\section{Introduction}\label{sec:intro}
Many mixed-integer models keep most of their constraints in a set $X$
and couple blocks, stages or scenarios through equalities $r(x)=0$.
Examples are linking constraints between subsystems, coupling
constraints in market models, copy constraints between a stage and its
state, and nonanticipativity constraints. We call the constraints that
define $X$ the \emph{native constraints} and the equalities $r(x)=0$
the \emph{relaxed} (or \emph{linking}) \emph{equalities}. The vector $x$ collects the integer and
continuous variables. The problem is
\[
 v=\min\{f(x):x\in X,\ r(x)=0\},
\]
and moving the relaxed equalities into the objective with a multiplier
$\lambda\in\R^m$ and a norm penalty with coefficient $\rho\ge0$ gives
\[
 L_\rho(\lambda)=\min_{x\in X}\{f(x)+\lambda^Tr(x)+\rho\norm{r(x)}\},
 \qquad
 D_\rho=\sup_{\lambda\in\R^m}L_\rho(\lambda).
\]
Section~\ref{sec:geometry} states the setting precisely
(compact $X$, continuous $f$ and $r$, a fixed norm). Weak duality gives
$L_\rho(\lambda)\le D_\rho\le v$. For $\rho=0$, $D_0$ is the ordinary
Lagrangian dual, which for mixed-integer problems can leave a gap
$D_0<v$. The norm term, which gives the sharp augmented Lagrangian (not
the squared term of the classical augmented Lagrangian), can close this
gap with a finite coefficient. Such a coefficient is an
\emph{exact penalty}.

Exact penalty functions were introduced by \citet{Eremin1966} and
\citet{Zangwill1967}, and \citet{Burke1991} unified their theory.
\citet{Rockafellar1974} showed that augmented Lagrangians remove duality
gaps in nonconvex programming. \citet{RockafellarWets1998} and
\citet{HuangYang2003} developed the connection between augmented
Lagrangian duality and exact penalization. For mixed-integer problems,
finite exact penalties are known for bounded pure integer programs,
rational MILPs, convex MIQPs, convex objectives over polyhedral
mixed-integer sets, and mixed-integer nonlinear problems under
regularity conditions on the integer slices
\citep{BolandEberhard2015,FeizollahiAhmedSun2017,GuAhmedDey2020,Bhardwaj2024,LefebvreSchmidt2025};
Section~\ref{sec:related} states these results precisely.

We use the following terms; Definition~\ref{def:exactness} states them
formally. A coefficient $\rho$ is \emph{value-exact} if $D_\rho=v$, that
is, if the penalized dual with the best multiplier has the optimal value.
It is \emph{exact at $\lambda$} if $L_\rho(\lambda)=v$ for that fixed
multiplier. It is \emph{minimizer-exact at $\lambda$} if the minimizers of
$f+\lambda^Tr+\rho\norm r$ over $X$ are exactly the optimal solutions of
the original problem. Minimizer-exactness at $\lambda$ implies exactness
at $\lambda$, which implies value-exactness; Section~\ref{sec:geometry}
shows that the converses fail. The \emph{least penalty}
$\rho_*=\inf\{\rho\ge0:D_\rho=v\}$ is, when finite, the smallest
value-exact coefficient, with the multiplier optimized. The
\emph{zero-multiplier threshold} $\rho_*(0)$ is, when finite, the
smallest coefficient that is exact at $\lambda=0$. We call a coefficient
\emph{sufficient} if it is value-exact, or exact at $\lambda=0$ when the
zero-multiplier threshold is meant.

\paragraph{Why the size of the coefficient matters.}
Exact norm penalties are used in decomposition methods, and each of them
must choose the coefficient. In cutting-plane methods for multistage stochastic MILPs,
augmented Lagrangian cuts with a norm as augmenting function give
Lipschitz lower approximations of cost-to-go functions; they are tight
for a suitable multiplier and large enough $\rho$, but arbitrarily
large $\rho$ may be needed near discontinuities of the value function
\citep{AhmedCabralCosta2022}.
\citet{ZhangSun2022} penalize copy constraints between a state and its
local copy, assume that this penalty reformulation is exact (their
Assumption~2), and their iteration bounds depend on the penalty sizes. For Lipschitz
regularization in stochastic dual dynamic programming, which penalizes
such copy constraints by a norm, a recent survey calls the choice of a
sufficiently large coefficient an open challenge in practice
\citep{FullnerRebennack2025}.
In $\ell_1$ augmented Lagrangian and ADMM decompositions of two-block
MILPs, the penalty that guarantees an $\varepsilon$-accurate solution
grows like $1/\varepsilon$, and such a large penalty slows the method
\citep{SunSunYin2024}; in bundle methods for generalized augmented
Lagrangian duals, adjusting the penalty parameter was the delicate part
\citep{CordovaOliveiraSagastizabal2022}; and penalty alternating
direction heuristics increase weighted $\ell_1$ penalties on copy
constraints in an outer loop \citep{GeisslerMorsiScheweSchmidt2017}.
In $\ell_1$ augmented Lagrangian pricing of nonconvex electricity
markets, the coefficient is part of the price signal, and
\citet{WangHesamzadehZhaoKronqvist2026} seek a valid coefficient as close
to the least one as is computationally feasible. Large coefficients
also cost precision: floating-point MIP solvers suffer from wide ranges
of data magnitudes \citep{KlotzNewman2013}, and quantum annealers rescale
QUBO coefficients to a fixed range, so a large penalty compresses the
gaps between objective values \citep{KrpanPovhPucher2025}.

These uses raise two questions about a sufficient coefficient. We
measure coefficients in ordinary binary encoding: a rational number is
given by the binary representations of its numerator and its positive
denominator.
\begin{itemize}
\item \emph{Representation.} How many bits does a sufficient coefficient
 need, as a function of the encoding length of the instance?
\item \emph{Calibration.} How closely can the least penalty $\rho_*$ be
 approximated, from above, by an efficient algorithm?
\end{itemize}
Bit length and numerical size are different measures: a coefficient with
polynomially many bits may still be exponentially large as a number. A
coefficient whose bit length is superpolynomial cannot even be written
down in polynomial time.

Both questions are posed in the conclusion of
\citet{LefebvreSchmidt2025}. They note that an exact penalty of
polynomial size is known only for MIQPs and ask whether this extends to
classes such as quadratically constrained quadratic problems. They also
observe that the coefficient computed by their polynomial-time method is
``too large to be useful in practical applications'', and they ask
whether the smallest penalty that closes the duality gap can be computed
in polynomial time, conjecturing that it cannot. We answer the first
question negatively for convex quadratic native constraints when both
the continuous dimension and the number of nonlinear constraints grow,
and positively when either is fixed, under explicit variable bounds and
a Slater condition on the integer slices (Assumption~\ref{ass:model});
for nonconvex quadratic constraints the question remains open
(Section~\ref{sec:discussion}). We confirm the conjecture unless
$\mathrm P=\mathrm{NP}$, even for approximation within any polynomial
factor.

\subsection{Main results}
Table~\ref{tab:results} summarizes the results. Throughout, $N$ is the
encoding length of an instance, $n$ the number of continuous variables
and $m$ the number of relaxed equalities.

\begin{table}[!t]
\centering
\small
\renewcommand{\arraystretch}{1.15}
\begin{tabular}{@{}>{\raggedright\arraybackslash}p{0.12\linewidth}
 >{\raggedright\arraybackslash}p{0.35\linewidth}
 >{\raggedright\arraybackslash}p{0.33\linewidth}
 >{\raggedright\arraybackslash}p{0.12\linewidth}@{}}
\toprule
Result & Setting and hypotheses & Conclusion & Reference\\
\midrule
Lower bound &
One binary and $n=d+1$ continuous variables; linear objective;
$d-1$ convex quadratic native inequalities ($d\ge2$); coefficients from a
fixed finite set; sparse encoding length $O(d\log(d+1))$; points with slack $1/8$
on every integer slice &
$\rho_*=2^{2^d-1}$, $\rho_*(0)=2^{2^d}$; every value-exact rational
coefficient has at least $2^d$ numerator bits; a gap at most
$\epsilon<1$ needs $2^d-O_\epsilon(1)$ bits &
Thm.~\ref{thm:lower}, Cor.~\ref{cor:accuracy}\\
\addlinespace
Lower bound in $N$ &
Same chain with seed $a_1\ge2^{-t}$, $t$ of order $N$ &
least penalty with $\Omega(N2^n)$ bits &
Rem.~\ref{rem:seed}\\
\addlinespace
Upper bound in $n$ &
Assumption~\ref{ass:model}: degree-two data, convex in the continuous
variables for each integer assignment; finite bounds; nonempty feasible
set; refined Slater condition; affine relaxed equalities &
an integer $\rho$ with $(N+1)2^{O(n)}$ bits, minimizer-exact at
$\lambda=0$ for both $\norm\cdot_\infty$ and $\norm\cdot_1$ &
Thm.~\ref{thm:generalupper}\\
\addlinespace
Upper bound in $k$ &
Assumption~\ref{ass:model} with at most $k$ native inequalities nonlinear
in the continuous variables; any number of affine rows &
the same with $N^{O(k+1)}$ bits; for fixed $n$ or fixed $k$ such a
coefficient can be printed in polynomial time &
Thm.~\ref{thm:fixedcount}, Cor.~\ref{cor:printing}\\
\addlinespace
Calibration &
Binary box, one equality, linear objective; known unique optimizer;
$\rho=1$ is exact at $\lambda=0$ &
unless $\mathrm P=\mathrm{NP}$, no polynomial-time algorithm finds
$\widehat\rho$ with $\rho_*\le\widehat\rho\le p(N)\rho_*$, and likewise
for $\rho_*(0)$; testing value-exactness of $\rho=1/3$ is
coNP-complete &
Thm.~\ref{thm:polyfactor}, Thm.~\ref{thm:coNP}\\
\addlinespace
Strong hardness &
Stable-set native set; coefficients in $\{-1,0,1\}$ &
$\rho_*=\alpha(G)$, so computing $\rho_*$ is strongly NP-hard &
Prop.~\ref{prop:stableset}\\
\addlinespace
Perturbation &
At most $K$ compact convex slices; continuous convex objectives with
supplied range $M$; affine residuals; right-hand side sampled from a
rational grid with at least $\lceil4mK/\epsilon\rceil$ points per
coordinate in a cube of radius $\sigma$ &
with probability at least $1-\epsilon$, the sampled problem is infeasible
or $4mKM/(\sigma\epsilon)$ is exact at $\lambda=0$ for
$\norm\cdot_\infty$; the order $K/\epsilon$ is necessary &
Thm.~\ref{thm:gridsmoothing}, Sec.~\ref{sec:numerical-size}\\
\bottomrule
\end{tabular}
\caption{Informal summary of the main results. $\rho_*$: least penalty
(multiplier optimized); $\rho_*(0)$: zero-multiplier threshold;
$p$: any fixed polynomial.}
\label{tab:results}
\end{table}

\paragraph{A doubly exponential lower bound.}
Theorem~\ref{thm:lower} gives, for each chain length $d$, an instance with
one binary variable, $n=d+1$ continuous variables, a linear objective,
one relaxed equality, and convex quadratic and affine native constraints
with coefficients from a fixed finite set; its sparse encoding length is
$O(d\log(d+1))$. Its least penalty is $2^{2^d-1}$, even after optimizing
the multiplier, so every exact rational coefficient, in any of the three
senses, needs at least $2^d$ numerator bits. This holds although every
integer slice has points with slack at least $1/8$, and a fixed additive
accuracy below the objective range does not help
(Corollary~\ref{cor:accuracy}).

The mechanism is simple. A chain of convex quadratic inequalities
$a_{i+1}\ge a_i^2$ lets the improving integer assignment come within
$\delta_d=2^{-2^d}$ of satisfying the relaxed equality, but never reach
it. The other assignment has residuals of both signs at zero cost. A
convex combination of its point with residual $-1$ and the improving
point, weighted so that the residuals average to zero (a \emph{balanced
mixture}, Section~\ref{sec:geometry}), cancels every multiplier, and its
average penalized cost is negative unless $\rho\ge1/(2\delta_d)$. The numbers grow fast. For $d=5$, with six
continuous variables, the improving assignment violates the relaxed
equality by $\delta_5=2^{-32}\approx2.3\cdot10^{-10}$, below common
feasibility tolerances, and the least penalty is $2^{31}$. For $d=10$ it
is $2^{1023}$, the largest power of two in IEEE double precision; for
$d=11$ it exceeds the double-precision range. Thus finite exactness with
convex quadratic native constraints does not come with the polynomial
encoding guarantee known for linear native constraints. The obstruction
lies in the representation of the penalty, not in the original problem,
whose optimal value is immediate.

\paragraph{Encoding upper bounds.}
For the positive results we fix the integer variables and look at the
resulting \emph{slices} of the native set. A slice is
\emph{equality-feasible} if it contains a point satisfying the relaxed
equalities. The \emph{refined Slater condition} requires, for every
equality-feasible slice, a point of the slice that satisfies the relaxed
equalities and strictly satisfies every native inequality that is
nonlinear in the continuous variables; inequalities affine in the
continuous variables may hold with equality, and nothing is assumed on
the other slices. Assumption~\ref{ass:model} asks for degree-two data
that are convex in the continuous variables for each integer
assignment, explicit finite bounds on all variables, affine relaxed
equalities, a nonempty feasible set and the refined Slater condition.
Under this assumption, Theorem~\ref{thm:generalupper} gives an integer
coefficient with $(N+1)2^{O(n)}$ bits, and
Theorem~\ref{thm:fixedcount} gives one with $N^{O(k+1)}$ bits when at
most $k$ native inequalities are nonlinear in the continuous variables;
the number of affine rows is arbitrary. In both cases one coefficient
is minimizer-exact at multiplier zero for the infinity norm and the
one-norm. By Remark~\ref{rem:seed}, the worst case of the first bound
is $N2^{\Theta(n)}$ bits.

Both proofs use Lemma~\ref{lem:slices}, which reduces exactness at
multiplier zero to a multiplier bound on each equality-feasible slice
and a lower bound on the residual of each equality-infeasible slice;
effective bounds from real algebraic geometry control both uniformly
over all integer assignments. Consequently, within the class of
Assumption~\ref{ass:model}, coefficients with superpolynomially many
bits require both a growing continuous dimension and a growing number
of nonlinear native inequalities; the lower-bound instance has both.
For fixed $n$, and separately for fixed $k$, a conservative coefficient
can be printed in polynomial time (Corollary~\ref{cor:printing}). These
are existence bounds: the coefficient may be numerically huge, the
constants are not evaluated, and nothing is said about solving the
penalized subproblem.

\paragraph{Hardness of calibration.}
Theorem~\ref{thm:polyfactor} considers pure binary linear programs with
one equality: the native set is a binary box and the objective is
linear. On these instances the unique optimal solution is known and
$\rho=1$ is exact at multiplier zero. Even so, unless
$\mathrm P=\mathrm{NP}$, no polynomial-time algorithm can output a
sufficient coefficient within any polynomial factor of the least
penalty, or of the zero-multiplier threshold; Section~\ref{sec:calibration}
extends this to factors $2^{N^{1-\epsilon}}$ for every fixed
$\epsilon\in(0,1)$. Deciding whether the
coefficient $\rho=1/3$ is value-exact is coNP-complete
(Theorem~\ref{thm:coNP}). The instances encode \textsc{Subset Sum}: when the target sum is
attainable, some integer assignment that lowers the objective has a
residual of magnitude $1$; otherwise every such assignment has a residual
of magnitude at least $\kappa-1$, for a scaling integer $\kappa$ that we
take exponentially large. The hardness is weak, since a pseudo-polynomial
dynamic program computes $\rho_*$; a stable-set variant with
coefficients in $\{-1,0,1\}$ makes exact computation strongly NP-hard
(Proposition~\ref{prop:stableset}). The message is that a safe
coefficient can be easy to find while a nearly least one is hard to
certify.

\paragraph{Right-hand-side perturbations.}
Theorem~\ref{thm:gridsmoothing} considers at most $K$ compact convex
slices with continuous convex objectives of range $M$ and affine
residual maps. If the right-hand side is sampled from a rational grid
with at least $\lceil4mK/\epsilon\rceil$ points per coordinate in a cube
of radius $\sigma$, then with probability at least $1-\epsilon$ the
sampled problem is infeasible or the infinity-norm coefficient
$4mKM/(\sigma\epsilon)$ is exact at multiplier zero. If the number of grid points per coordinate is the smallest power of
two that is at least $\lceil4mK/\epsilon\rceil$, the grid points and the
coefficient have encoding length polynomial in the supplied rational
data, in $m$ and in $\log K$ (Remark~\ref{rem:gridencoding}). The mechanism is that
large penalties arise only when the right-hand side lies near the
boundary of the set of residuals attainable on some slice; a random
right-hand side avoids these boundaries with high probability, and a
convex repair argument turns the distance into a coefficient. The
guarantee concerns the perturbed problem, whose optimal integer
assignment may differ from the original one. The coefficient is
numerically of order $K$, which can be exponential in the number of
integer variables, and Section~\ref{sec:numerical-size} shows that the
order $K/\epsilon$ is necessary and that the expected least penalty can
be infinite.

\subsection{Related work}\label{sec:related}
\paragraph{Exact penalties and augmented Lagrangian duality.}
The dual of \citet{Rockafellar1974} maximizes jointly over the
multipliers and the penalty parameter; we fix $\rho$, optimize
$\lambda$, and ask for the least $\rho$ at which the optimized value is
exact.\footnote{Theorem, section and example numbers of preprints and
author manuscripts refer to the versions identified in the references.}
Section~\ref{sec:geometry} relates the threshold $\rho_*(\lambda)$ to
the exact penalty representations of the sharp Lagrangian in
\citet[Chapter~11]{RockafellarWets1998}. In the study of global
saddle points of augmented Lagrangians for continuous cone-constrained
problems, \citet{Dolgopolik2018} uses the term ``least exact penalty
parameter''; our least penalty $\rho_*$ is the analogous quantity for
mixed-integer problems, taken after optimizing the multiplier.

\paragraph{Exactness and encoding for mixed-integer problems.}
\citet{BolandEberhard2015} prove finite-penalty exactness for bounded
pure integer programs under their assumptions on the augmenting
function, and \citet{FeizollahiAhmedSun2017} prove norm-penalty exactness
for rational MILPs. \citet[Theorem~11]{GuAhmedDey2020} show that for a
convex quadratic objective over a rational mixed-integer polyhedron some
exact norm penalty has polynomial encoding length.
\citet[Theorems~3.1--3.3]{Bhardwaj2024} give value-function proofs of
existence with data-dependent penalty parameters for rational MILPs and
for MIQPs with bounded integer variables. For convex objectives over
polyhedral mixed-integer sets they prove exactness when the integer
variables are bounded on the feasible set or when the objective and the
continuous relaxation share no nonzero recession direction, and they
give an explicit bound under strong convexity and smoothness.
\citet[Theorem~14]{LefebvreSchmidt2025} prove that a finite penalty is
exact at each fixed multiplier under compactness, a differentiable
convex objective, convex slices and a Slater condition on every
equality-feasible slice; the instance of Theorem~\ref{thm:lower}
satisfies these assumptions. Their Theorem~15 computes, in polynomial
time, an infinity-norm penalty that is exact at a given rational
multiplier for MILPs with integer data, and their Table~1 remarks that
the extension to convex MIQPs is easy. For native constraints linear in
all variables and a jointly convex quadratic objective, polynomial
encoding length of an exact penalty is therefore prior work
\citep[Theorem~11]{GuAhmedDey2020}, and so is polynomial-time output of
an exact coefficient for MILPs with integer data
\citep[Theorem~15]{LefebvreSchmidt2025}. Appendix~\ref{app:attainment}
discusses their Example~13: under their supremum-based Definition~1 it
is value-exact at every penalty, but no multiplier attains the dual
value.

\paragraph{Squaring chains and bit complexity.}
The squaring chain and its small residuals are classical;
\citet[Example~6.1]{Bienstock2022} attribute it to Ramana and Khachiyan.
Example~6.2 of the same paper, a nonconvex QCQP with a point of doubly
exponentially small violation, already implies large zero-multiplier
norm penalties. \citet[Sections~2 and~4, Result~4]{Beck2023} use a
bounded convex squaring chain with Slater regularity to show that bilevel
problems are sensitive to feasibility errors. To our knowledge, what is
new in Theorem~\ref{thm:lower} is the exact threshold of the optimized dual for
a compact convex-quadratic model with one binary variable and uniform
strict points, together with its fixed-accuracy consequence. Residuals
of opposite signs on the cost-free slice make the bound hold for every
multiplier.

\paragraph{Real-algebraic tools.}
The upper bounds rest on established tools: radii of balls meeting the
connected components of semialgebraic sets \citep{BasuRoy2010} and
quantitative quantifier elimination \citep{BasuPollackRoy2006,Basu2014}.
Algebraic bounds that depend on the number of quadratic constraints
appear in \citet{GrigorievPasechnik2005} and \citet{KammingaRudolph2025};
Section~\ref{sec:fixedcount} compares them with our argument. To our
knowledge, what is new here are the resulting encoding bounds for exact
penalties, including the dependence on the number of native inequalities
that are nonlinear in the continuous variables while the number of
affine rows is arbitrary.

\paragraph{Hardness of choosing penalties.}
Hardness of choosing penalties, even when the constrained optimum is
known, predates this paper. \citet[Section~IV.A, Lemma~1]{Alessandroni2025}
prove hardness of choosing penalties for exact QUBO reformulations. Their
reduction imposes $x=0$ with a squared residual and a quadratic
objective; on binary variables their penalty $M\sum_ix_i$ is reproduced
by a linear multiplier with augmentation coefficient zero, so it does not
give a positive optimized threshold. \citet[Section~2.2]{Mirkarimi2024}
note that a universally successful, polynomially tuned linear Ising
penalty for knapsack would imply $\mathrm P=\mathrm{NP}$, allowing that no
linear penalty may work on some instances. On the family of
Theorem~\ref{thm:polyfactor}, $D_0=-1<v$, so
no purely linear (Lagrangian) penalty is ever exact and the norm term is
essential. \citet[Definition~1 and Theorem~2]{Alessandroni2026} study
probabilistic feasibility and objective guarantees when exact Gibbs
sampling is used for approximate optimization. Safe conservative
penalties are also known: \citet[Sections~5--6]{GusmeroliWiegele2022}
obtain exact squared penalties for binary quadratic problems with
integer-data linear equality constraints from objective bounds and
residual separation ($\norm{Ax-b}^2\ge1$ by integrality). To our
knowledge, what is new in Theorem~\ref{thm:polyfactor} is the polynomial-factor conclusion for a
norm-penalty threshold that is always positive and is taken after
unrestricted optimization of the multiplier, under the restrictions
above: a binary box, one equality, a linear objective, a known unique
optimizer and a known sufficient coefficient.

\paragraph{Conditioning and perturbation.}
Distance to the boundaries of convex sets is a classical conditioning
device \citep{DunaganSpielmanTeng2011};
\citet[Theorem~3.3 and Corollary~3.4]{BuergisserAmelunxen2012} bound
the relative volume of boundary neighbourhoods of spherically convex
sets within spherical caps.
Theorem~\ref{thm:gridsmoothing} is an explicit application to exact
penalties, with a rational grid, an encoding bound and the infeasibility
alternative; its Gaussian variant (Proposition~\ref{prop:gaussian}) has
the infeasibility alternative but no grid or encoding statement. Their ingredients are known: the boundary-avoidance
principle, the convex repair argument and, for the Gaussian variant, the
tube estimate \citep[Lemma~2.3.4]{DunaganSpielmanTeng2011}. The relation between convex exact
penalties and multiplier bounds is classical
\citep[Theorem~4.2]{FriedlanderTseng2007}, and Lipschitz value bounds
across discrete choices appear in mixed-integer optimal control
\citep[Theorem~3]{GugatHante2017}.

\paragraph{Organization.}
Section~\ref{sec:geometry} defines the exactness notions and the two
thresholds and proves the balanced-mixture, one-equality and slice
results. Section~\ref{sec:lower} proves the lower bound and discusses
which features of the formulation it depends on.
Sections~\ref{sec:upper} and~\ref{sec:fixedcount} prove the encoding
upper bounds in the continuous dimension and in the number of nonlinear
native inequalities. Section~\ref{sec:calibration} proves the
calibration hardness, and Section~\ref{sec:perturbation} the perturbation
results. Section~\ref{sec:discussion} discusses consequences, scope and
open questions. Appendix~\ref{app:attainment} treats Example~13 of
\citet{LefebvreSchmidt2025}, and Appendix~\ref{app:symmetric} gives a
symmetric two-binary variant of the lower bound in which the optimized
and zero-multiplier thresholds coincide.
